% GENERATED FILE. Edit canonical lessons, then run npm run generate:books.
% canonical-source-hash: fnv1a64:86c92371f8633e02
% canonical-lessons: JA-C492-nieru, JA-C492-waku, JA-C492-hitasu, JA-C492-fukasu, JA-C492-muku

\chapter{Simmer, Soak and Steam}
\label{ch:ja-simmer-soak-and-steam}
\hlchaptermodality{492}

\begin{chapteropening}
\textbf{By the end of this chapter:} \emph{I can say how food is cooked.}

Complete the last lesson of chapter 492: I can say how food is cooked.
\end{chapteropening}

\section[\texorpdfstring{\ja{にえる} (nieru)}{にえる (nieru)}]{\ja{にえる} (nieru) — to be cooked through}

\label{lesson:JA-C492-nieru}



\begin{quote}
Before the new one: say the Japanese for tap water, then the Japanese for to come to the boil.
\end{quote}

\subsection*{You'll want to know: \ja{にえる}}
\textbf{\ja{にえる}} — \emph{nieru} — \textquotedblleft{}to be cooked through\textquotedblright{}.

\begin{quote}
\textbf{\ja{だいこんが} \ja{よく} \ja{にえました。}} — \emph{daikon ga yoku niemashita} — \textquotedblleft{}The radish is well cooked.\textquotedblright{}
\end{quote}

Polite \textbf{\ja{にえます}}, \emph{niemasu}; an ichidan verb. Past (polite): \textbf{\ja{だいこんが} \ja{にえました。}}, \emph{daikon ga niemashita}.

\begin{cousinweb}
\textbf{\ja{にえる}} (\emph{nieru}) is a native Japanese word (wago), inherited rather than borrowed from Chinese; its kanji was attached later as a kun'yomi, and its older origin is not settled (JMdict, KANJIDIC2).
\end{cousinweb}

\begin{grammarlens}[title={What you've built}]
One more word for describing what you lived and felt.
\end{grammarlens}

\subsection*{Your turn}
\begin{itemize}
\raggedright
  \item \emph{Say it:} \emph{nieru}
  \item \emph{Say it:} \emph{nieru}, once more
  \item \emph{Say it:} \emph{daikon ga yoku niemashita} — The radish is well cooked.
  \item \emph{Recall:} say the Japanese for tap water, then the Japanese for to come to the boil, then say \emph{nieru} again
\end{itemize}

\begin{tcolorbox}[breakable,colback=teal!4,colframe=teal!35!black,title={Before you move on}]
What does \ja{にえる} mean? (\textquotedblleft{}To be cooked through\textquotedblright{}.) Say it once more.
\end{tcolorbox}
\section[\texorpdfstring{\ja{わく} (waku)}{わく (waku)}]{\ja{わく} (waku) — to boil; to get hot}

\label{lesson:JA-C492-waku}



\begin{quote}
Before the new one: say the Japanese for to come to the boil, then the Japanese for to be cooked through.
\end{quote}

\subsection*{You'll want to know: \ja{わく}}
\textbf{\ja{わく}} — \emph{waku} — \textquotedblleft{}to boil; to get hot\textquotedblright{}.

\begin{quote}
\textbf{\ja{やかんの} \ja{ゆが} \ja{わいています。}} — \emph{yakan no yu ga waiteimasu} — \textquotedblleft{}The water in the kettle is boiling.\textquotedblright{}
\end{quote}

Polite \textbf{\ja{わきます}}, \emph{wakimasu}; a godan verb. Past (polite): \textbf{\ja{おふろが} \ja{わきました。}}, \emph{ofuro ga wakimashita}.

\begin{cousinweb}
\textbf{\ja{わく}} (\emph{waku}) is a native Japanese word (wago), inherited rather than borrowed from Chinese; its kanji was attached later as a kun'yomi, and its older origin is not settled (JMdict, KANJIDIC2).
\end{cousinweb}

\begin{grammarlens}[title={What you've built}]
One more word for describing what you lived and felt.
\end{grammarlens}

\subsection*{Your turn}
\begin{itemize}
\raggedright
  \item \emph{Say it:} \emph{waku}
  \item \emph{Say it:} \emph{waku}, once more
  \item \emph{Say it:} \emph{yakan no yu ga waiteimasu} — The water in the kettle is boiling.
  \item \emph{Recall:} say the Japanese for to come to the boil, then the Japanese for to be cooked through, then say \emph{waku} again
\end{itemize}

\begin{tcolorbox}[breakable,colback=teal!4,colframe=teal!35!black,title={Before you move on}]
What does \ja{わく} mean? (\textquotedblleft{}To boil; to get hot\textquotedblright{}.) Say it once more.
\end{tcolorbox}
\section[\texorpdfstring{\ja{ひたす} (hitasu)}{ひたす (hitasu)}]{\ja{ひたす} (hitasu) — to dip in liquid; to steep}

\label{lesson:JA-C492-hitasu}



\begin{quote}
Before the new one: say the Japanese for to be cooked through, then the Japanese for to boil.
\end{quote}

\subsection*{You'll want to know: \ja{ひたす}}
\textbf{\ja{ひたす}} — \emph{hitasu} — \textquotedblleft{}to dip in liquid; to steep\textquotedblright{}.

\begin{quote}
\textbf{\ja{まめを} \ja{みずに} \ja{ひたします。}} — \emph{mame o mizu ni hitashimasu} — \textquotedblleft{}I soak the beans in water.\textquotedblright{}
\end{quote}

Polite \textbf{\ja{ひたします}}, \emph{hitashimasu}; a godan verb. Past (polite): \textbf{\ja{ははが} \ja{まめを} \ja{みずに} \ja{ひたしました。}}, \emph{haha ga mame o mizu ni hitashimashita}.

\begin{cousinweb}
\textbf{\ja{ひたす}} (\emph{hitasu}) is a native Japanese word (wago), inherited rather than borrowed from Chinese; its kanji was attached later as a kun'yomi, and its older origin is not settled (JMdict, KANJIDIC2).
\end{cousinweb}

\begin{grammarlens}[title={What you've built}]
One more word for describing what you lived and felt.
\end{grammarlens}

\subsection*{Your turn}
\begin{itemize}
\raggedright
  \item \emph{Say it:} \emph{hitasu}
  \item \emph{Say it:} \emph{hitasu}, once more
  \item \emph{Say it:} \emph{mame o mizu ni hitashimasu} — I soak the beans in water.
  \item \emph{Recall:} say the Japanese for to be cooked through, then the Japanese for to boil, then say \emph{hitasu} again
\end{itemize}

\begin{tcolorbox}[breakable,colback=teal!4,colframe=teal!35!black,title={Before you move on}]
What does \ja{ひたす} mean? (\textquotedblleft{}To dip in liquid; to steep\textquotedblright{}.) Say it once more.
\end{tcolorbox}
\section[\texorpdfstring{\ja{ふかす} (fukasu)}{ふかす (fukasu)}]{\ja{ふかす} (fukasu) — to steam (food)}

\label{lesson:JA-C492-fukasu}



\begin{quote}
Before the new one: say the Japanese for to boil, then the Japanese for to dip in liquid.
\end{quote}

\subsection*{You'll want to know: \ja{ふかす}}
\textbf{\ja{ふかす}} — \emph{fukasu} — \textquotedblleft{}to steam (food)\textquotedblright{}.

\begin{quote}
\textbf{\ja{いもを} \ja{ふかして} \ja{たべます。}} — \emph{imo o fukashite tabemasu} — \textquotedblleft{}I steam potatoes and eat them.\textquotedblright{}
\end{quote}

Polite \textbf{\ja{ふかします}}, \emph{fukashimasu}; a godan verb. Past (polite): \textbf{\ja{そぼが} \ja{いもを} \ja{ふかしました。}}, \emph{sobo ga imo o fukashimashita}.

\begin{cousinweb}
\textbf{\ja{ふかす}} (\emph{fukasu}) is a native Japanese word (wago), inherited rather than borrowed from Chinese; its kanji was attached later as a kun'yomi, and its older origin is not settled (JMdict, KANJIDIC2).
\end{cousinweb}

\begin{grammarlens}[title={What you've built}]
One more word for describing what you lived and felt.
\end{grammarlens}

\subsection*{Your turn}
\begin{itemize}
\raggedright
  \item \emph{Say it:} \emph{fukasu}
  \item \emph{Say it:} \emph{fukasu}, once more
  \item \emph{Say it:} \emph{imo o fukashite tabemasu} — I steam potatoes and eat them.
  \item \emph{Recall:} say the Japanese for to boil, then the Japanese for to dip in liquid, then say \emph{fukasu} again
\end{itemize}

\begin{tcolorbox}[breakable,colback=teal!4,colframe=teal!35!black,title={Before you move on}]
What does \ja{ふかす} mean? (\textquotedblleft{}To steam (food)\textquotedblright{}.) Say it once more.
\end{tcolorbox}
\section[\texorpdfstring{\ja{むく} (muku)}{むく (muku)}]{\ja{むく} (muku) — to peel}

\label{lesson:JA-C492-muku}



\begin{quote}
Before the new one: say the Japanese for to dip in liquid, then the Japanese for to steam (food).
\end{quote}

\subsection*{You'll want to know: \ja{むく}}
\textbf{\ja{むく}} — \emph{muku} — \textquotedblleft{}to peel\textquotedblright{}.

\begin{quote}
\textbf{\ja{りんごの} \ja{かわを} \ja{むきます。}} — \emph{ringo no kawa o mukimasu} — \textquotedblleft{}I peel the apple.\textquotedblright{}
\end{quote}

Polite \textbf{\ja{むきます}}, \emph{mukimasu}; a godan verb. Past (polite): \textbf{\ja{あねが} \ja{みかんを} \ja{むきました。}}, \emph{ane ga mikan o mukimashita}.

\begin{cousinweb}
\textbf{\ja{むく}} (\emph{muku}) is a native Japanese word (wago), inherited rather than borrowed from Chinese, and usually written in kana; its older origin is not settled (JMdict).
\end{cousinweb}

\begin{grammarlens}[title={What you've built}]
One more word for describing what you lived and felt.
\end{grammarlens}

\subsection*{Your turn}
\begin{itemize}
\raggedright
  \item \emph{Say it:} \emph{muku}
  \item \emph{Say it:} \emph{muku}, once more
  \item \emph{Say it:} \emph{ringo no kawa o mukimasu} — I peel the apple.
  \item \emph{Recall:} say the Japanese for to dip in liquid, then the Japanese for to steam (food), then say \emph{muku} again
\end{itemize}

\begin{tcolorbox}[breakable,colback=teal!4,colframe=teal!35!black,title={Before you move on}]
What does \ja{むく} mean? (\textquotedblleft{}To peel\textquotedblright{}.) Say it once more.
\end{tcolorbox}
